% !PRINTED 402
\weberpage{402}
%<<CONT
\[
x_1 = a_1 + \frac{1}{x_2}
\]
in $f_2(x_2)$. Die Gleichung $f_2 = 0$ hat dann wieder mindestens
eine Wurzel, die grösser als 1 ist, und die also zwischen den
beiden positiven ganzen Zahlen $a_2$ und $a_2 + 1$ liegt, u.\,s.\,f.
Schliesslich wird man auf diese Weise dazu kommen, eine einzige
Wurzel zu isoliren, und man erhält in der Reihe der Näherungsbrüche
\[
(a_0), \quad (a_0, a_1), \quad (a_0, a_1, a_2), \quad (a_0, a_1, a_2, a_3) \ldots
\]
eine Reihe rationaler Zahlen, die abwechselnd über und unter
einer der Wurzeln $x$ liegen und sich dieser Wurzel unbegrenzt
annähern. Ist man so bis zu dem Näherungsbruch
\[
(a_0, a_1, a_2 \ldots a_{\nu-1}) = \frac{P_\nu}{Q_\nu}
\]
vorgedrungen, so unterscheidet sich dieser von dem wahren Werth
von $x$ um weniger als $1 : Q_\nu^2$. Die Grösse des Nenners giebt
also unmittelbar ein Maass für die bereits erreichte Genauigkeit.

In den gewöhnlichen Fällen wird man die Zahlen $a_0, a_1, a_2 \ldots$
einfach dadurch erhalten, dass man auf die Zeichenwechsel der
Functionen $f, f_1, f_2 \ldots$ achtet.

Um dies Verfahren an einem Beispiel zu erproben, nehmen
wir die cubische Gleichung, die wir schon früher betrachtet
haben:
\[
x^3 - 2x - 2 = 0.
\]
Sie hat eine Wurzel zwischen 1 und 2, also haben wir $a_0 = 1$
zu setzen.

Man erhält, wenn man die oben angegebene Transformation
wirklich ausführt, die folgende Kette von Gleichungen:
\[
\begin{aligned}
x^3 - 2x - 2 &= 0, & a_0 &= 1\\
3x^3 - x^2 - 3x - 1 &= 0, & a_1 &= 1\\
2x^3 - 4x^2 - 8x - 3 &= 0, & a_2 &= 3\\
9x^3 - 22x^2 - 14x - 2 &= 0, & a_4 &= 1\\ % [sic] printed order a_4 before a_3
46x^3 - 6x^2 - 32x - 9 &= 0, & a_3 &= 2\\
x^3 - 94x^2 - 132x - 46 &= 0, & a_5 &= 95\\
3561x^3 - 9183x^2 - 191x - 1 &= 0, & a_6 &= 2,
\end{aligned}
\]
also
\[
x = (1, 1, 3, 2, 1, 95, 2 \ldots)
\]
und die Näherungsbrüche:
%CONT>>
